\SourceLeaf{292}
$(A,0,C)$ ad semissem reducere possumus, retinendo e binis semper unam; et in
omnibus remanentibus erit $A<\sqrt{\pm D}$.

Simili modo si inter formas art. praec. occurrit forma $(2b,b,c)$, inter easdem reperietur
\[
(4c-2b,\ 2c-b,\ c)=\left(-\frac{2D}{b},\ -\frac{D}{b},\ c\right)
\]
quae illi proprie aequivalens et ab ipsa diversa erit, unico quem seponimus casu
excepto, ubi $c=b=\pm1$ sive $D=-1$. Ex his duabus formis eam retinere
sufficit, cuius terminus primus est minor quam terminus primus alterius (magnitudine
aequales, signis diversi in hoc casu esse nequeunt); unde patet, etiam
omnes formas $(2B,B,C)$ ad semissem reduci posse, e binis unam semper eiiciendo;
et in remanentibus esse $B<\frac{D}{B}$ sive $B<\sqrt{\pm D}$. Hoc modo ex omnibus
formis art. praec. semissis tantum remanet, quarum complexum per $W$ designabimus,
nihilque superest, nisi ut ostendamus, quot classes diversae ex his
formis oriantur. Ceterum manifestum est, in eo casu ubi $D$ sit negativus, totidem
formas positivas in $W$ affore quot negativas.

I. Quando $D$ est negativus, singulae formae in $W$ pertinebunt ad classes
diversas. Nam omnes formae $(A,0,C)$ erunt reductae; similiter omnes formae
$(2B,B,C)$ reductae erunt, praeter eas in quibus $C<2B$; in tali vero forma
erit $2C<2B+C$; unde (quoniam $B<\frac{D}{B}$ \textit{i. e.}\ $B<2C-B$, adeoque $2B<2C$,
sive $B<C$), $2C-2B<C$ et $C-B<\frac12 C$ et proin $(C,C-B,C)$, quae
manifesto illi aequivalet, forma reducta. Hoc modo totidem formae reductae
habentur, quot formae habentur in $W$, et quum facile perspiciatur, inter illas
neque identicas neque oppositas occurrere posse, (unico casu excepto, ubi
$C-B=0$, in quo erit $B=C=\pm1$, adeoque $D=-1$, quem iam seposuimus):
omnes ad classes diversas pertinebunt. Hinc colligitur, multitudinem
omnium classium ancipitum pr. primitivarum det. $D$ multitudini formarum in
$W$ seu semissi multitudinis formarum art. praec. aequalem esse; in casu excepto
autem $D=-1$ per compensationem idem evenit, scilicet duae classes
habentur, ad quarum alteram pertinent formae $(1,0,1)$, $(2,1,1)$, ad alteram hae
$(-1,0,-1)$, $(-2,-1,-1)$. Generaliter itaque pro determinante negativo
multitudo omnium classium ancipitum pr. prim. aequalis est multitudini omnium
\SourceLeaf{293}
characterum assignabilium formarum primitivarum huius determinantis; multitudo
classium ancipitum pr. prim. positivarum autem semissis erit.

II. Quando $D$ est positivus quadratus $=hh$, haud difficile demonstratur,
singulas formas in $W$ ad classes diversas pertinere; sed pro hoc casu ad problematis
solutionem adhuc brevius sequenti modo pervenire possumus. Quum per
art. 210 in quavis classe ancipite pr. prim. det. $hh$, neque in ulla alia, contineatur
forma reducta una $(a,h,0)$, in qua $a$ est valor expr. $\sqrt{1}\ (\mathrm{mod.}\ 2h)$ inter $0$ et
$2h-1$ incl. situs: perspicuum est, totidem classes ancipites pr. prim. det. $hh$
dari, quot valores expressio illa habeat. Ex art. 105 autem nullo negotio deducitur,
multitudinem horum valorum esse $2^n$ vel $2^{n+1}$ vel $2^{n+2}$, prout $h$ sit
impar vel impariter par vel pariter par, sive prout $D\equiv1$ vel $\equiv4$ vel $\equiv0\ (\mathrm{mod.}\ 8)$,
designante $n$ multitudinem divisorum primorum imparium ipsius $h$ sive ipsius $D$.
Hinc colligitur, multitudinem classium ancipitum pr. prim. semper esse semissem
multitudinis omnium formarum in art. praec. erutarum, sive multitudini formarum
in $W$ vel omnium characterum possibilium aequalem.

III. Quando $D$ est positivus non-quadratus, ex singulis formis $(A,B,C)$
in $W$ contentis alias deducamus $(A,B',C')$, accipiendo $B'\equiv B\ (\mathrm{mod.}\ A)$ et inter
limites $\sqrt{D}$ et $\sqrt{D}\mp A$ (ubi signum superius vel inferius adhibendum, prout $A$.
est pos. vel neg.) atque $C'=\frac{B'B'-D}{A}$; designemusque harum complexum per $W'$.
Manifesto hae formae erunt proprie primitivae ancipites det. $D$, atque omnes
inter se diversae: praeterea vero omnes erunt formae reductae. Quando enim
$A<\sqrt{D}$, $B'$ manifesto erit $<\sqrt{D}$ atque positivus; praeterea $B'>\sqrt{D}\mp A$
adeoque $A>\sqrt{D}-B'$ et proin $A$, positive acceptus, certo inter $\sqrt{D}+B'$ et
$\sqrt{D}-B'$ situs. Quando vero $A>\sqrt{D}$, non poterit esse $B=0$ (quippe quas
formas eiecimus), sed erit necessario $B=\frac12 A$; hinc $B'$ magnitudine ipsi $\frac12 A$
aequalis, signo positivus (quoniam enim $A<2\sqrt{D}$, $\pm\frac12 A$ iacebit inter limites
ipsi $B'$ assignatos, ipsique $B$ sec. mod. $A$ erit congruus; quare $B'=\pm\frac12 A$),
proin $B'<\sqrt{D}$, unde $2B'<\sqrt{D}+B'$ sive $A<\sqrt{D}+B'$, quamobrem $\pm A$
necessario inter limites $\sqrt{D}+B'$ et $\sqrt{D}-B'$ iacebit. Denique $W'$ omnes formas
reductas pr. prim. ancipites det. $D$ continebit; si enim $(a,b,c)$ est huiusmodi
forma, erit vel $b\equiv0$, vel $b\equiv\frac12 a\ (\mathrm{mod.}\ a)$. In casu priori manifesto non poterit
esse $b<a$ neque adeo $a>\sqrt{D}$, quapropter forma $\left(a,0,-\frac{D}{a}\right)$ certo contenta
\SourceLeaf{294}
erit in $W$, et respondens $(a,b,c)$ in $W'$; in posteriori certo erit $a<2\sqrt{D}$,
adeoque $\left(a,\allowbreak \frac12 a,\allowbreak \frac14 a-\frac{D}{a}\right)$ in $W$ contenta, atque respondens $(a,b,c)$ in $W'$.
Ex his colligitur, multitudinem formarum in $W$ aequalem esse multitudini
omnium formarum reductarum ancipitum pr. prim. det. $D$; quoniam vero in singulis
classibus ancipitibus \textit{binae} formae reductae ancipites continentur (artt. 187,
194), multitudo omnium classium ancipitum pr. prim. det. $D$ erit semissis multitudinis
formarum in $W$, sive semissis multitudinis omnium characterum assignabilium.

\Needspace{8\baselineskip}
\Article{259}
Multitudo classium ancipitum improprie primitivarum determinantis dati $D$
multitudini proprie primitivarum eiusdem det. semper est aequalis. Sit $K$ classis
principalis, atque $K'$, $K''$ etc. reliquae classes ancipites pr. primitivae huius determinantis;
$L$ aliqua classis anceps improprie primitiva eiusdem det., \textit{e. g.}\ ea
in qua est forma $(2,1,\frac12-\frac12 D)$. Prodibit itaque ex compositione classis $L$ cum $K$
classis $L$ ipsa; ex compositione classis $L$ cum $K'$, $K''$ etc. provenire supponamus
classes $L'$, $L''$ etc. resp., quae manifesto omnes ad eundem determinantem
$D$ pertinebunt, atque improprie primitivae et ancipites erunt. Patet itaque,
theorema demonstratum fore, simulac probatum fuerit, omnes classes $L$, $L'$, $L''$
etc. esse diversas, aliasque ancipites impr. prim. det. $D$ praeter illas non dari.
Ad hunc finem sequentes casus distinguimus:

I. Quando multitudo classium impr. primitivarum multitudini pr. primitivarum
aequalis est, quaevis illarum oritur ex compositione classis $L$ cum classe
determinata proprie primitiva, unde necessario omnes $L$, $L'$, $L''$ etc. erunt diversae.
Designante autem $\mathfrak{L}$ classem quamcunque ancipitem impr. prim. det. $D$,
dabitur classis proprie primitiva $\mathfrak{K}$ talis ut sit $\mathfrak{K}+L=\mathfrak{L}$; si classi $\mathfrak{K}$ opposita
est classis $\mathfrak{K}'$, erit etiam (quoniam classes $L$, $\mathfrak{L}$ sibi ipsae oppositae sunt)
$\mathfrak{K}'+L=\mathfrak{L}$, unde necessario $\mathfrak{K}$ cum $\mathfrak{K}'$ identica erit, adeoque classis anceps:
hinc $\mathfrak{K}$ reperietur inter classes $K$, $K'$, $K''$ etc. atque $\mathfrak{L}$ inter has $L$, $L'$, $L''$ etc.

II. Quando multitudo classium improprie primitivarum ter minor est quam
multitudo classium pr. primitivarum, sit $H$ classis in qua est forma $\left(4,\allowbreak 1,\allowbreak \frac{1-D}{4}\right)$,
$H'$ ea in qua est forma $\left(4,\allowbreak 3,\allowbreak \frac{9-D}{4}\right)$, eruntque $H$, $H'$ proprie primitivae et tum
\SourceLeaf{295}
inter se tum a classe principali $K$ diversae, atque $H+H'=K$; $2H=H'$,
$2H'=H$; et si $\mathfrak{L}$ est classis quaecunque improprie primitiva det. $D$, quae oritur
ex compositione classis $L$ cum proprie primitiva $\mathfrak{K}$, erit etiam $\mathfrak{L}=L+\mathfrak{K}+H$
et $\mathfrak{L}=L+\mathfrak{K}+H'$; praeter tres classes (pr. prim. atque diversas) $\mathfrak{K}$, $\mathfrak{K}+H$,
$\mathfrak{K}+H'$ aliae non dabuntur, quae cum $L$ compositae ipsam $\mathfrak{L}$ producant. Quoniam
igitur, si $\mathfrak{L}$ est anceps atque $\mathfrak{K}'$ ipsi $\mathfrak{K}$ opposita, etiam $L+\mathfrak{K}'=\mathfrak{L}$, necessario
$\mathfrak{K}'$ cum aliqua illarum trium classium identica erit. Si $\mathfrak{K}'=\mathfrak{K}$, erit $\mathfrak{K}$ anceps;
si $\mathfrak{K}'=\mathfrak{K}+H$, erit $K=\mathfrak{K}+\mathfrak{K}'=2\mathfrak{K}+H=2(\mathfrak{K}+H')$ adeoque $\mathfrak{K}+H'$
anceps; similiterque si $\mathfrak{K}'=\mathfrak{K}+H'$ erit $\mathfrak{K}+H$ anceps, unde concluditur,
$\mathfrak{L}$ inter classes $L$, $L'$, $L''$ etc. necessario reperiri. Facile autem perspicitur, inter
tres classes $\mathfrak{K}$, $\mathfrak{K}+H$, $\mathfrak{K}+H'$ plures ancipites esse non posse; si enim tum $\mathfrak{K}$
tum $\mathfrak{K}+H$ ancipites essent sive cum oppositis suis $\mathfrak{K}'$, $\mathfrak{K}'+H'$ resp. identicae,
foret $\mathfrak{K}+H=\mathfrak{K}+H'$; eadem conclusio resultat ex suppositione, $\mathfrak{K}$ et $\mathfrak{K}+H'$
esse ancipites; denique si $\mathfrak{K}+H$, $\mathfrak{K}+H'$ ancipites sive cum oppositis suis
$\mathfrak{K}'+H'$, $\mathfrak{K}'+H$ identicae essent, fieret $\mathfrak{K}+H+\mathfrak{K}'+H=\mathfrak{K}'+H'+\mathfrak{K}+H'$,
unda $2H=2H'$, sive $H'=H$. Quamobrem unica tantum classis anceps pr.
prim. dabitur, quae cum $L$ composita ipsam $\mathfrak{L}$ producit, adeoque omnes $L$, $L'$, $L''$
etc. erunt diversae.

Multitudo classium ancipitum in ordine \textit{derivato} manifesto aequalis est multitudini
classium ancipitum in ordine primitivo, ex quo est derivatus, adeoque per
praecedentia semper poterit assignari.

\Needspace{8\baselineskip}
\Article{260}
\textsc{Problema.} \textit{Classis proprie primitiva $K$ determinantis $D$ oritur ex duplicatione
classis proprie primitivae $k$ eiusdem determinantis: quaeruntur o\,m\,n\,e\,s similes
classes, ex quarum duplicatione classis $K$ oritur.}

\textit{Sol.} Sit $H$ classis principalis det. $D$ atque $H'$, $H''$, $H'''$ etc. reliquae
classes ancipites pr. primitivae eiusdem determinantis; classes quae ex harum
compositione cum $k$ oriuntur, $k+H'$, $k+H''$, $k+H'''$ designentur per $k'$, $k''$, $k'''$
etc. Tunc omnes classes $k$, $k'$, $k''$ etc. erunt pr. primitivae det. $D$ et inter se
diversae; aeque facile perspicitur, ex singularum duplicatione oriri classem $K$.
Denotante autem $\mathfrak{K}$ classem quamcunque pr. prim. det. $D$, quae duplicata producit
classem $K$, necessario inter classes $k$, $k'$, $k''$ etc. contenta erit. Ponatur
enim $\mathfrak{K}=k+\mathfrak{H}$, ita ut $\mathfrak{H}$ sit classis pr. prim. det. $D$ (art. 249). eritque
\SourceLeaf{296}
$2k+2\mathfrak{H}=2\mathfrak{K}=K=2k$, unde facile concluditur, $2\mathfrak{H}$ coincidere cum classe
principali, $\mathfrak{H}$ esse ancipitem sive inter $H$, $H'$, $H''$ etc. contentam, atque $\mathfrak{K}$ inter
$k$, $k'$, $k''$ etc.; quamobrem hae classes completam problematis solutionem exhibent.

Ceterum manifestum est, in eo casu, ubi $D$ sit negativus, e classibus
$k$, $k'$, $k''$ etc. semissem fore classes positivas, semissem negativas.

Quum igitur quaevis classis pr. prim. det. $D$, quae ex ullius classis similis
duplicatione oriri potest, omnino ex totidem classium similium duplicatione proveniat,
quot classes ancipites pr. prim. det. $D$ dantur: perspicuum est, si multitudo
cunctarum classium pr. prim. det. $D$ sit $r$, multitudo omnium classium
ancipitum pr. prim. huius det. $n$, multitudinem omnium classium pr. prim. eiusdem
det. quae ex duplicatione similis classis produci possint, fore $\frac{r}{n}$. Eadem
formula resultat, si, pro det. negativo, characteres $r$, $n$ multitudinem classium
\textit{positivarum} designant, ille \textit{omnium} pr. prim., hic solarum ancipitum. Ita \textit{e. g.} pro
$D=-161$ multitudo omnium classium pr. prim. positivarum est $16$, multitudo
ancipitum $4$, unde multitudo omnium classium, quae per duplicationem alicuius
classis oriri possunt, debebit esse $4$. Et revera invenitur, omnes classes in genere
principali contentas hac proprietate esse praeditas; scilicet classis principalis
$(1,0,161)$ oritur ex duplicatione quatuor classium ancipitum; $(2,1,81)$ ex duplicatione
classium $(9,1,18)$, $(9,-1,18)$, $(11,2,15)$, $(11,-2,15)$; $(9,1,18)$
ex dupl. classium $(3,1,54)$, $(6,1,27)$, $(5,-2,33)$, $(10,3,17)$; denique
$(9,-1,18)$ ex duplicatione classium $(3,-1,54)$, $(6,-1,27)$, $(5,2,33)$,
$(10,-3,17)$.

\Needspace{11\baselineskip}
\begin{center}
\textit{Certe semissi omnium characterum pro determinante dato assignabilium genera proprie primitiva
(positiva pro det. neg.) respondere nequeunt.}
\end{center}

\Needspace{8\baselineskip}
\Article{261}
\textsc{Theorema.} \textit{Semissi omnium characterum assignabilium pro determinante positivo
non-quadrato nulla genera proprie primitiva respondere possunt; pro determinante
negativo autem nulla genera proprie primitiva positiva.}

\textit{Dem.} Sit $m$ multitudo omnium generum proprie primitivorum (positivorum)
determinantis $D$; $k$ multitudo classium in singulis generibus contentarum,
ita ut $km$ sit multitudo omnium classium proprie primitivarum (positivarum);
$n$ multitudo omnium characterum diversorum pro hoc det. assignabilium. Tunc
\SourceLeaf{297}
per art. 258 multitudo omnium classium ancipitum (positivarum) pr. primitivarum
erit $\frac12 n$; hinc per art. praec. multitudo omnium classium pr. prim., quae ex duplicatione
similis classis oriri possunt, erit $\frac{2km}{n}$. Sed per art. 247 hae classes omnes
pertinent ad genus principale, in quo continentur $k$ classes; si itaque omnes
classes generis principalis ex duplicatione alicuius classis provenire possunt (quod
revera semper locum habere in sequentibus demonstrabitur), erit $\frac{2km}{n}=k$, sive
$m=\frac12 n$; certo autem nequit esse $\frac{2km}{n}>k$ neque adeo $m>\frac12 n$. Quoniam itaque
multitudo omnium generum pr. prim. (positivorum) certo non est maior quam
semissis omnium characterum assignabilium: ad minimum horum semissi talia
genera respondere nequeunt. \textit{Q. E. D.} --- Ceterum probe notandum est, hinc
nondum sequi, semissi omnium characterum assignabilium revera respondere genera
pr. prim. (positiva), sed huius propositionis gravissimae veritas infra demum
e reconditissimis numerorum mysteriis enodari poterit.

Quum pro determinante negativo totidem genera negativa semper exstent
quot positiva, manifesto ex omnibus characteribus assignabilibus non plures quam
semissis generibus pr. prim. negativis competere possunt, de qua re ut et de generibus
impr. prim. infra loquemur. Denique observamus, theorema ad determinantes
positivos quadratos non extendi, pro quibus nullo negotio perspicitur
singulis characteribus assignabilibus genera revera respondere.

\Needspace{11\baselineskip}
\begin{center}
\textit{Theorematis fundamentalis et reliquorum theorematum ad residua $-1$, $+2$, $-2$ pertinentium
demonstratio secunda.}
\end{center}

\Needspace{8\baselineskip}
\Article{262}
In eo itaque casu, ubi pro determinante non-quadrato dato $D$ duo tantummodo
characteres diversi assignari possunt, unico tantum genus pr. primitivum
(positivum) respondebit, (quod non poterit esse aliud quam genus principale),
alter nulli formae pr. prim. (pos.) illius determinantis competet. Hoc evenit pro
determinantibus $-1$, $2$, $-2$, $-4$. numeris primis formae $4n+1$ positive, iisque
formae $4n+3$ negative acceptis, denique pro omnibus numerorum primorum
formae $4n+1$ potestatibus exponentis imparis positive sumtis, et pro potestatibus
numerorum primorum formae $4n+3$ positive vel negative sumtis prout
exponentes sunt pares vel impares. Ex hoc principio methodum novam haurire
possumus, non modo theorema fundamentale, sed etiam reliqua theoremata Sect.
praec. ad residua $-1$, $+2$, $-2$ pertinentia demonstrandi, quae a methodis in
\SourceLeaf{298}
Sect. praec. adhibitis omnino est diversa, elegantiaque his neutiquam inferior
aestimanda videtur. Determinantem $-4$ autem, et qui sunt numerorum primorum
potestates, quum nihil novi doceant, praeteribimus.

Pro determinante $-1$ itaque nulla forma positiva datur, cuius character sit
$3,4$; pro determinante $+2$ nulla omnino forma, cuius character sit $3\ \mathit{et}\ 5,8$;
pro determinante $-2$ nulli formae positivae competet character $5\ \mathit{et}\ 7,8$; pro
determinante $+p$, si $p$ est numerus primus formae $4n+1$, vel pro determinante
$-p$, si $p$ est numerus primus formae $4n+3$, nulli formae pr. pr. (positivae
in casu post.) competet character $Np$. Hinc theoremata Sect. praec. sequenti
modo demonstramus:

I. Est $-1$ non-residuum cuiusvis numeri (positivi) formae $4n+3$. Si
enim $-1$ residuum talis numeri $A$ esset, faciendo $-1=BB-AC$, foret
$(A,B,C)$ forma positiva det. $-1$, cuius character $3,4$.

II. Est $-1$ residuum cuiusvis numeri primi $p$ formae $4n+1$. Nam
character formae $(-1,0,p)$, sicuti omnium proprie primitivarum det. $p$. erit
$Rp$, adeoque $-1Rp$.

III. Tum $+2$ tum $-2$ est residuum cuiusvis numeri primi $p$ formae
$8n+1$. Nam vel formae $\left(8,\allowbreak 1,\allowbreak \frac{1-p}{8}\right)$, $\left(-8,\allowbreak 1,\allowbreak \frac{p-1}{8}\right)$, vel hae $\left(8,\allowbreak 3,\allowbreak \frac{9-p}{8}\right)$, $\left(-8,\allowbreak 3,\allowbreak \frac{p-9}{8}\right)$
erunt proprie primitivae (prout $n$ impar vel par), adeoque ipsarum character $Rp$;
hinc $+8Rp$ et $-8Rp$, unde etiam $2Rp$, $-2Rp$.

IV. Est $+2$ non-residuum cuiusvis numeri formae $8n+3$ aut $8n+5$.
Si enim esset residuum talis numeri $A$, daretur forma $(A,B,C)$ determinantis
$+2$, cuius character $3\ \mathit{et}\ 5,8$.

V. Simili modo $-2$ est non-residuum cuiusvis numeri formae $8n+5$
aut $8n+7$, alioquin enim daretur forma $(A,B,C)$ determinantis $-2$, cuius
character $5\ \mathit{et}\ 7,8$.

VI. Est $-2$ residuum cuiusvis numeri primi $p$ formae $8n+3$. Hanc
propositionem per methodum duplicem demonstrare licet. \textit{Primo}, quum per IV
sit $+2Np$, atque per I, $-1Np$, necessario erit $-2Rp$. Demonstratio \textit{secunda}
petitur ex consideratione determinantis $+2p$, pro quo quatuor characteres
sunt assignabiles, puta $Rp$, $1\ \mathit{et}\ 3,8$; $Rp$, $5\ \mathit{et}\ 7,8$; $Np$, $1\ \mathit{et}\ 3,8$; $Np$, $5\ \mathit{et}\ 7,8$,
ex quibus igitur saltem duobus nulla genera respondebunt. Iam formae
\SourceLeaf{299}
$(1,0,-2p)$ competit character primus; formae $(-1,0,2p)$ quartus; quare qui
reiici debent sunt secundus atque tertius. Quum itaque character formae
$(p,0,-2)$ relative ad numerum $8$ sit $1\ \mathit{et}\ 3,8$, ipsius character relative ad $p$ non
poterit esse alius quam $Rp$, unde $-2Rp$.

VII. Est $+2$ residuum cuiusvis numeri primi $p$ formae $8n+7$ quod
per methodum duplicem demonstrare licet. \textit{Primo}, quum ex I et V sit $-1Np$,
$-2Np$, erit $+2Rp$. \textit{Secundo} quum vel $\left(8,\allowbreak 1,\allowbreak \frac{1+p}{8}\right)$ vel $\left(8,\allowbreak 3,\allowbreak \frac{9+p}{8}\right)$ sit forma
proprie primitiva determinantis $-p$ (prout $n$ par vel impar), ipsius character
erit $Rp$, adeoque $8Rp$ et $2Rp$.

VIII. Quilibet numerus primus $p$ formae $4n+1$ est non-residuum cuiusvis
numeri imparis $q$, qui ipsius $p$ non-residuum est. Patet enim, si $p$ esset
residuum ipsius $q$, dari formam proprie primitivam determinantis $p$, cuius character
$Np$.

IX. Simili modo si numerus quicunque impar $q$ est non-residuum numeri
primi $p$ formae $4n+3$, erit $-p$ non-residuum ipsius $q$; alioquin enim daretur
forma positiva pr. primitiva determinantis $-p$ cuius character $Np$.

X. Quivis numerus primus $p$ formae $4n+1$ est residuum cuiusvis alius
numeri primi $q$, qui ipsius $p$ residuum est. Si etiam $q$ est formae $4n+1$, hoc
statim sequitur ex VIII; si vero $q$ est formae $4n+3$, erit etiam $-q$ residuum
ipsius $p$ (propter II) adeoque $pRq$ (ex IX).

XI. Si numerus quicunque primus $q$ est residuum alius numeri primi $p$
formae $4n+3$, erit $-p$ residuum ipsius $q$. Si enim $q$ est formae $4n+1$; ex
VIII sequitur $pRq$, adeoque (per II), $-pRq$; casus autem ubi etiam $q$ est formae
$4n+3$, huic methodo se subducit, attamen facile ex consideratione determinantis
$+pq$ absolvi potest. Scilicet quum ex quatuor characteribus pro hoc
determinante assignabilibus $Rp$. $Rq$; $Rp$, $Nq$; $Np$, $Rq$; $Np$, $Nq$ duobus nulla
genera respondere possint, atque formarum $(1,0,-pq)$, $(-1,0,pq)$ characteres
respective sint primus et quartus, character secundus et tertius nulli formae
pr. prim. det. $pq$ competere possunt. Quum itaque character formae $(q,0,-p)$
resp. numeri $p$ per hyp. sit $Rp$, eiusdem formae character respectu numeri $q$
debet esse $Rq$, adeoque $-pRq$. \textit{Q. E. D.}

Si in proposs. VIII et IX, $q$ supponitur designare numerum primum, hae
cum X et XI iunctae theorema fundamentale Sect. praec. exhibent.
\Needspace{11\baselineskip}
\SourceLeaf{300}
\begin{center}
\textit{Ea characterum semissis, quibus genera respondere nequeunt, propius determinatur.}
\end{center}

\Needspace{8\baselineskip}
\Article{263}
Postquam theorema fundamentale demonstratione nova comprobavimus, eam
characterum semissem, quibus nullae formae pr. primitivae (positivae) respondere
possunt, pro determinante quocunque non-quadrato dato discernere ostendemus,
quod negotium eo brevius absolvere licebit, quum ipsius fundamentum iam in
disquisitione artt. 147--150 sit contentum. Sit $ee$ quadratum maximum, determinantem
propositum $D$ metiens, atque $D=D'ee$, ita ut $D'$ nullum factorem
quadratum implicet; porro sint $a$, $b$, $c$ etc. omnes divisores primi impares ipsius
$D'$, adeoque $D'$ sine respectu signi sui vel productum ex his numeris vel duplum
huius producti. Designetur per $\Omega$ complexus characterum particularium $Na$, $Nb$,
$Nc$ etc., solus, quando $D'\equiv1\ (\mathrm{mod.}\ 4)$; adiuncto charactere $3,4$, quando $D'\equiv3$
atque $e$ impar aut impariter par; adiunctis his $3,8$ atque $7,8$, quando $D'\equiv3$
atque $e$ pariter par; adiuncto vel charactere $3\ \mathit{et}\ 5,8$, vel duobus $3,8$ atque $5,8$,
quando $D'\equiv2\ (\mathrm{mod.}\ 8)$ atque $e$ vel impar vel par; denique adiuncto vel charactere
$5\ \mathit{et}\ 7,8$, vel duobus $5,8$ atque $7,8$, quando $D'\equiv6\ (\mathrm{mod.}\ 8)$ atque $e$
vel impar vel par. His ita factis, omnibus characteribus integris, in quibus
multitudo impar characterum particularium $\Omega$ continetur, nulla genera proprie
primitiva (positiva) determinantis $D$ respondere poterunt. In omnibus casibus
characteres particulares, qui exprimunt relationem ad tales divisores primos ipsius
$D$, qui ipsum $D'$ non metiuntur, ad generum possibilitatem vel impossibilitatem
nihil conferunt. --- Ex theoria combinationum autem facillime perspicitur,
hoc modo revera semissem omnium characterum integrorum assignabilium excludi.

Demonstratio horum praeceptorum adornatur sequenti modo. E principiis
Sect. praec., sive theorematibus in art. praec. denuo demonstratis nullo negotio
deducitur, si $p$ sit numerus primus (impar positivus) ipsum $D$ non metiens, cui
aliquis e characteribus reiectis competat, $D'$ implicare multitudinem imparem
factorum, qui sint non-residua ipsius $p$, atque adeo $D'$, et hinc etiam $D$, esse
non-residuum ipsius $p$; porro facile perspicitur, productum e numeris quotcunque
imparibus ad $D$ primis, quorum nulli aliquis characterum reiectorum competat,
etiam cum tali charactere consentire non posse; hinc vice versa perspicuum
est, quemvis numerum imparem positivum ad $D$ primum, cui aliquis characte%
\SourceLeaf{301}%
rum reiectorum conveniat, certe aliquem factorem primum eiusdem qualitatis implicare,
adeoque $D$ ipsius non-residuum esse. Si itaque forma proprie primitiva
(positiva) determinantis $D$ daretur, alicui characterum reiectorum respondens,
$D$ foret non-residuum cuiusvis numeri positivi imparis ad ipsum primi per talem
formam repraesentabilis, quod manifesto cum theoremate art. 154 consistere
nequit.

Tamquam exempla conferantur classificationes in artt. 231, 232 traditae,
quarum numerum quisque pro lubitu augere poterit.

\Needspace{8\baselineskip}
\Article{264}
Hoc itaque modo pro quovis determinante non-quadrato dato omnes characteres
assignabiles in duas species $P$, $Q$ aequaliter distribuuntur, ita ut nulli
characterum $Q$ forma proprie primitiva (positiva) respondere possit, reliquis autem
$P$, quantum quidem hucusque novimus, nihil obstet, quominus ad tales formas
pertineant. Circa has characterum species notetur imprimis propositio sequens,
quae ex ipsarum criterio facile deducitur: Si character ex $P$ cum charactere
ex $Q$ componitur (ad normam art. 246 perinde ac si etiam huic genus
responderet) prodibit character ex $Q$; si vero duo characteres ex $P$. vel duo ex
$Q$ componuntur, character resultans ad $P$ pertinebit. Adiumento huius theorematis
etiam pro generibus negativis atque improprie primitivis semissis omnium
characterum assignabilium excludi potest sequenti modo.

I. Pro determinante negativo $D$ genera negativa positivis hoc respectu
prorsus contraria erunt, scilicet nullus characterum $P$ pertinebit ad genus proprie
primitivum negativum, sed haec genera omnia habebunt characteres ex $Q$.
Quando enim $D'\equiv1\ (\mathrm{mod.}\ 4)$, erit $-D'$ numerus positivus formae $4n+3$,
adeoque inter $a$, $b$, $c$ etc. multitudo impar numerorum formae $4n+3$, quorum
singulorum non-residuum erit $-1$, unde patet, in characterem integrum formae
$(-1,0,D)$ in hoc casu ingredi multitudinem imparem characterum particularium
ex $\Omega$, sive illum pertinere ad $Q$; quando $D'\equiv3\ (\mathrm{mod.}\ 4)$, ex simili ratione
inter $a$, $b$, $c$ etc. vel nullus numerus formae $4n+3$ reperietur, vel duo, vel
quatuor etc., sed quum vel $3,4$ vel $3,8$ vel $7,8$ in hoc casu occurrat inter characteres
particulares formae $(-1,0,D)$, patet, characterem integrum huius formae
etiam hic\textperiodcentered{}pertinere ad $Q$. Eadem conclusio aeque facile in casibus reliquis
